% !TEX root = ../main.tex
% capitulos/3_resultados_discussao.tex

\subsection{Apresentação de figuras, tabelas e equações}

\textbf{Figuras} e \textbf{Tabelas} devem ser inseridas diretamente no corpo do texto, utilizando preferencialmente o mesmo padrão (tamanho de letra, borda, etc.). Quando citadas no texto, escrever com a 1ª letra maiúscula e não abreviar.

``Na Figura 1 é possível observar a evolução da população...''; ``... De acordo com a Tabela 2 ...''

As \textbf{Equações} quando citadas no texto virão com a 1ª letra maiúscula e o número entre parênteses, sem abreviação.

``Obtendo-se assim a Equação (1):''

As \textbf{figuras} deverão ser numeradas em algarismos arábicos, por ordem de aparição no texto e devem estar centralizadas. Para mais de uma \textbf{figura} em um mesmo contexto, utilizar \textbf{figuras} compostas no qual devem ser identificadas com a numeração sequencial e uma letra identificadora maiúscula (Figura 1A, Figura 1A-B). As \textbf{figuras} devem apresentar boa qualidade (resolução mínima de \textbf{300 dpi}) nos formatos TIFF ou JPEG. Sempre que necessários os autores devem providenciar adicionalmente os arquivos originais para editoração.

As legendas \textbf{devem ser concisas e autoexplicativas}, colocadas \textbf{abaixo} da \textbf{Figura}, com apenas a 1ª letra maiúscula na palavra ``\textbf{Figura}'' e no ``título'', sendo separado por dois pontos. A fonte usada para na legenda é a padrão usado em todo o texto (\textbf{Arial Narrow}), tamanho 11 e todo o texto da legenda deverá estar em Arial, 11.

\textbf{Figura 2}: para o caso de uma 2ª figura exposta no artigo.

Figura 2: Mapa viário do estado do Piauí, nordeste do Brasil, destacando as principais rodovias federais e estaduais.

\begin{figure}[htb]
    \centering
    \fbox{\rule{0pt}{2in} \rule{0.6\textwidth}{0pt}}
\end{figure}

Fonte: \textbf{Elaborado pelos autores, 2026.}

Quando houver mais de um gráfico para uma mesma figura, o título pode aparecer uma única vez, logo abaixo do conjunto de gráficos dispostos horizontal ou verticalmente.

\textbf{Figura 5}: para o caso de uma 5ª figura exposta no artigo.

Figura 5: Evolução de ganhos (em azul) e custos (em roxo) da produção da cajuína no estado do Piauí nos últimos 24 meses. A – histograma e gráfico de dispersão; B – apenas histograma.

\begin{figure}[htb]
    \centering
    \fbox{\rule{0pt}{2in} \rule{0.6\textwidth}{0pt}}
\end{figure}

Fonte: Dados da Pesquisa, 2024.

As \textbf{tabelas} deverão ser enumeradas em algarismos arábicos, por ordem de aparição no texto e devem estar centralizadas. O tamanho da fonte do texto interno da tabela é 11, sem espaçamento entre as linhas, o texto da primeira linha deverá vir em negrito, as bordas deverão seguir o padrão estabelecido no exemplo abaixo.

O título deverá vir \textbf{acima} com apenas a 1ª letra maiúscula na palavra ``Tabela'' e no ``título'', sendo separado por dois pontos. As unidades referentes à coluna, quando couber, serão apresentadas nos ``cabeçalhos'' da coluna correspondente. A fonte usada para no título da tabela é a padrão usado em todo o texto (\textbf{Arial Narrow}), o tamanho é 10 e apenas o título inicial deverá estar em negrito.

\textbf{Tabela 2}: para o caso de uma 2ª tabela exposta artigo

\textbf{Tabela 2}: Estudo da influência do tempo na degradação da glicose.

\begin{table}[htb]
    \centering
    \begin{tabular}{ccc}
        \toprule
        \textbf{Amostra} & \textbf{Concentração (moles/L)} & \textbf{Rendimento (\%)} \\
        \midrule
        1 & 0,02 & 45 \\
        2 & 0,12 & 56 \\
        3 & 0,30 & 70 \\
        4 & 0,43 & 87 \\
        \bottomrule
    \end{tabular}
    
    {\footnotesize Fonte: Dados da Pesquisa, 2024.}
\end{table}

As \textbf{equações} deverão estar enumeradas por ordem de aparição, com o respectivo número entre parênteses e no extremo da margem direita. Quando ocorrerem equações seguidas no texto, inserir uma linha como espaço entre as equações.

\begin{equation}
    E = mc^2
\end{equation}

\begin{equation}
    a^2 + b^2 = c^2
\end{equation}

\textbf{O sistema de unidades} deverá ser homogêneo em todo o texto. Recomenda-se o sistema internacional (SI).

Quanto ao \textbf{uso de palavras em língua estrangeira}, recomenda-se evitar o estrangeirismo. Quando o uso for necessário, utilizar a forma em itálico.

``O polímero produzido na etapa de finalização é estruturado na forma de \textit{pellets} ou \textit{flakes}''.
